\subsection{Original affine convexity and actual native regular-unit presentations}

\begin{definition}[The original affine native/Berkovich homeomorphism]
\label{berkovich:actual-affine-paper-homeomorphism}
\lean{EnvelopeProperty.Berkovich.RelativeSchemeValuedPoint.affineBerkHomeomorph}
\leanok
For every algebra $A$ over a field $k$, the native points of the actual
$\operatorname{Spec} A$ over the trivial valuation of $k$ are canonically
homeomorphic to the original affine Berkovich space of $A$.
\end{definition}
\begin{proof}
Compose the actual affine native-point/relative-seminorm homeomorphism
with the inverse of the original trivial-field Berkovich/relative-seminorm
homeomorphism. Both factors retain every original ring value.
\end{proof}

\begin{lemma}[The whole inverse affine point]
\label{berkovich:actual-affine-paper-homeomorphism-inverse}
\lean{EnvelopeProperty.Berkovich.RelativeSchemeValuedPoint.affineBerkHomeomorph_symm_coe}
\leanok
\uses{berkovich:actual-affine-paper-homeomorphism}
The inverse homeomorphism sends an original seminorm $z$ to its canonical
actual affine native point, including the original residue-field valuation.
\end{lemma}
\begin{proof}
The two relative subtypes impose the same trivial coefficient condition.
After forgetting that proof, their inverse homeomorphism is definitionally
the canonical affine spectrum/native-point inverse.
\end{proof}

\begin{lemma}[Original convexity gives actual regular units]
\label{berkovich:actual-affine-convexity-from-paper}
\lean{EnvelopeProperty.Berkovich.RelativeSchemeValuedPoint.IsConvexPLOn.of_paperAffine}
\leanok
\uses{berkovich:actual-affine-paper-homeomorphism,paper:original-affine-local-psh-comparison}
Let $U$ be any original affine analytic locus and $u$ any extended-real
function. If $u$ is convex PL on $U$ in the original paper definition,
its pullback is convex PL in actual scheme regular units on the whole
inverse-image native locus. No presentation-preservation premise is needed.
\end{lemma}
\begin{proof}
The proven original/relative-spectrum comparison converts the original
finite-coefficient presentations to relative-spectrum presentations, with
the same neighborhoods and nonvanishing coefficients. The genuine affine
regular-unit comparison replaces the coefficients by their actual regular
sections on a principal open where their product is invertible. Their
values, every rational affine expression, and its finite maximum are
unchanged. The two inverse homeomorphisms cancel on the original locus.
\end{proof}

\begin{lemma}[Actual regular units give original convexity]
\label{berkovich:actual-affine-convexity-to-paper}
\lean{EnvelopeProperty.Berkovich.RelativeSchemeValuedPoint.IsConvexPLOn.to_paperAffine}
\leanok
\uses{berkovich:actual-affine-paper-homeomorphism}
Conversely, every genuine intrinsic regular-unit convex presentation on an
actual affine native locus gives an original convex PL presentation after
the inverse homeomorphism, including its entire analytic neighborhood.
\end{lemma}
\begin{proof}
The actual affine comparison writes each original regular unit, near the
given point, as a fraction of actual elements of $A$ whose numerator and
denominator have nonzero value. Split its rational logarithmic term into
numerator and denominator terms. This produces the actual relative-spectrum
finite-coefficient presentation. Apply the original/relative-spectrum
comparison, retaining the whole finite maxima and neighborhoods. This is
the claimed pullback under the composite affine homeomorphism.
\end{proof}
